\documentclass[10pt,a4paper]{amsart}
\usepackage[margin=19mm]{geometry}
\usepackage{amsmath,amssymb,mathtools}
\usepackage[scr=rsfs,cal=euler]{mathalfa}
\usepackage{xfrac}
\usepackage[unicode,hidelinks,bookmarks=false]{hyperref}
\newcommand{\ie}{i.e.\ }
\newcommand{\cf}{cf.\ }
\newcommand{\resp}{respectively\ }
\newcommand{\Subsection}{\subsection}
\providecommand{\nobreakdash}{\nobreak\hbox{-}}
\setlength{\emergencystretch}{2em}
\multlinegap0pt
\usepackage{bm}
\usepackage{bbm}
\usepackage{dsfont}
\usepackage{xspace}
\usepackage{cancel}
\usepackage{mathtools}
\usepackage{amsthm}
\makeatletter
\@ifundefined{thm}{\newtheorem{thm}{Theorem}}{}
\makeatother
\newcommand{\Ac}{\mathcal{A}}
\newcommand{\RR}{\mathbb{R}}
\newcommand{\ZZ}{\mathbb{Z}}
\title[A new class of magic positive Ehrhart polynomials of reflexi]{A new class of magic positive Ehrhart polynomials of reflexive polytopes}
\author{Masato Konoike}
\date{Source: arXiv:2409.16648v2 (CC BY 4.0)}
\begin{document}
\maketitle

\noindent\emph{Source and licence.} A new class of magic positive Ehrhart polynomials of reflexive polytopes. Version arXiv:2409.16648v2 (CC BY 4.0), distributed under the Creative Commons Attribution 4.0 International licence, as recorded in the arXiv licence metadata for this exact version and shown on the abstract page https://arxiv.org/abs/2409.16648v2. Full rights notice: licenses/arxiv-2409.16648-CC-BY-4.0.txt.\par\smallskip

\noindent\textit{Editorial scope.} Counted unit: the Thm whose source label is \texttt{None} in the article (source0.tex lines 227--232, source label \texttt{None}), delivered together with the article's own complete original proof of it (source0.tex lines 233--302, one \texttt{proof} environment of 70 lines). No mathematics is written, repaired or re-derived: statement and proof are the article's own text line for line. Required context retained uncounted: none. Every definition, display and label of those lines is retained unchanged. Omitted from this package and not redistributed: 1--226, 303--511. Nothing inside a retained excerpt is omitted or altered: every retained body is delivered exactly as the article's lines are written, so no omission receipt of any kind accompanies this package. Citations: the retained excerpts carry no citation command at all, so no bibliography entry of the article is transcribed and no reference is redistributed. Reference adaptation: none was needed and none was made; every reference inside the delivered unit resolves inside it, so no unresolved reference or citation is produced. Printed slips kept literal: no printed word, symbol, hypothesis or number of the counted statement or of its proof was altered. Layout and class: the article's own class is \texttt{amsart} with packages including graphicx, xcolor, hyperref, amssymb, amstext, amsmath, amscd, amsthm, amsfonts, enumerate, caption, color, here, bm; the wrapper is the lane's pinned \texttt{amsart} editorial wrapper at 10pt on A4, which restates the theorem-like environments the retained text needs and adds the article's own macro definitions that the retained text needs (\texttt{\textbackslash Ac}, \texttt{\textbackslash RR}, \texttt{\textbackslash ZZ}), with extra wrapper packages (bm, bbm, dsfont, xspace, cancel, mathtools). The delivered numbering is the unit's own. Assets: no retained span contains a figure environment, a graphics-inclusion command, a PDF-page inclusion, a table or a TikZ picture; no image file is copied into this package, the package declares no asset dependency and no image file is redistributed with it. Licence: version arXiv:2409.16648v2 of this article is distributed under the Creative Commons Attribution 4.0 International licence, as recorded in the arXiv licence metadata for this exact version and shown on the abstract page https://arxiv.org/abs/2409.16648v2. A full rights notice is delivered at licenses/arxiv-2409.16648-CC-BY-4.0.txt. Characters: every retained excerpt is pure ASCII, so the delivered engine, XeLaTeX with its default Latin Modern text font, reports no missing character.
\begin{thm}
For $d \geq 2, E_{\mathrm{St}_d}(n)$ satisfies the following recurrence:
\begin{align}
  E_{\mathrm{St}_d^\ast}(n) = (2n+1)E_{\mathrm{St}_{d-1}^\ast}(n) - \frac{1}{2}n(n+1)E_{\mathrm{St}_{d-2}^\ast}(n).
\end{align}
\end{thm}

\begin{proof}
First, we define some sets. By the definition of the Stasheff polytope,
\[
\mathrm{St}_d^\ast = \{(x_1, \ldots, x_d) \in \RR^d \mid -1 \leq x_i \leq 1, x_j + \cdots + x_k \leq 1 (1 \leq i \leq d, 1 \leq j < k \leq d)\},
\]
\[
\mathrm{St}_{d-1}^\ast \times [-1, 1] = \{(x_1, \ldots, x_d) \in \RR^d \mid -1 \leq x_i \leq 1, x_j + \cdots + x_k \leq 1 (1 \leq i \leq d, 1 \leq j < k \leq d-1)\}.
\]
Let
$$\Ac = (\mathrm{St}_{d-1}^\ast \times [-1, 1]) \backslash \mathrm{St}_d^\ast.$$
For integers $m$ with $-n \leq m \leq n$, let
\[
A_m = \left(\bigcup_{i = 1}^{d - 2}(n\mathrm{St}_{d-2}^\ast \cap \{(x_1, \ldots, x_{d-2}) \in \RR^{d-2} \mid x_i + \cdots + x_{d-2} = m\})\right) \backslash \bigcup_{i = m+1}^{n}A_i,
\]
\[
\Delta_m = \{(x_{d-1}, x_d) \in \RR^2 \mid x_{d-1} \leq n-\max\{m, 0\}, x_d \leq n, x_{d-1} + x_d > n -\max\{m, 0\}\}.
\]
The following equation can be easily verified:
\begin{align}
  E_{\mathrm{St}_{d-1}^\ast \times [-1, 1]}(n) = (2n+1)E_{\mathrm{St}_{d-1}^\ast}(n).
\end{align}
Since $\mathrm{St}_d^\ast \subset \mathrm{St}_{d-1}^\ast \times [-1, 1]$, by (3.2)
\begin{align}
  \left|n\Ac \cap \ZZ^d\right| = (2n+1)E_{\mathrm{St}_{d-1}^\ast}(n) - E_{\mathrm{St}_d^\ast}(n).
\end{align}
Since $\Delta_m$ and $\{(x,y) \in \RR^2 \mid x \geq 0, y \geq 0, x+y < n\}$ are unimodularly equivalent, we have the following equation:
\begin{align}
  |\Delta_m \cap \ZZ^2| = \frac{1}{2}n(n+1).
\end{align}

\smallskip

\noindent
\underline{The first step}: We prove the following equality:
\begin{align}
  \bigsqcup_{-n \leq m \leq n}A_m \cap \ZZ^{d-2} = n\mathrm{St}_{d-2}^\ast \cap \ZZ^{d-2}.
\end{align}
Since $\bigsqcup_{-n \leq m \leq n}A_m \cap \ZZ^{d-2} \subset n\mathrm{St}_{d-2}^\ast \cap \ZZ^{d-2}$ is clear, we will prove $\bigsqcup_{-n \leq m \leq n}A_m \cap \ZZ^{d-2} \supset n\mathrm{St}_{d-2}^\ast \cap \ZZ^{d-2}$. For $x = (x_1, \ldots, x_{d-2}) \in n\mathrm{St}_{d-2}^\ast \cap \ZZ^{d-2}$, we define the value $p$ as follows:
$$p = \max{\{x_i + \cdots + x_{d-2} \mid 1 \leq i \leq d-2\}}.$$
Then, since $-n \leq p \leq n$ holds, and $x \in A_p \cap \ZZ^{d-2}$ follows from the definition of $A_m$. Therefore, we get (3.5).

From (3.4) and (3.5), we have the following equation:
\begin{align}
  \left|\bigsqcup_{-n \leq m \leq n}(A_m \times \Delta_m) \cap \ZZ^d \right| = \frac{1}{2}n(n+1)E_{\mathrm{St}_{d-2}^\ast}(n).
\end{align}

\smallskip

\noindent
\underline{The second step}: Next, we prove the following equality:
\begin{align}
  n\Ac \cap \ZZ^d = \bigsqcup_{-n \leq m \leq n}(A_m \times \Delta_m) \cap \ZZ^d.
\end{align}
Let $x = (x_1, \ldots, x_d) \in n\Ac \cap \ZZ^d$. Then, $x_d \leq n$ holds, and since $(x_1, \ldots, x_{d-2}) \in n\mathrm{St}_{d-2}^\ast \cap \ZZ^{d-2}$, that is, $(x_1, \ldots, x_{d-2}) \in \bigsqcup_{-n \leq m \leq n}A_m \cap \ZZ^{d-2}$ holds by (3.5), there exists an integer $-n \leq k \leq n$ such that $(x_1, \ldots, x_{d-2}) \in A_k$. 

Now, suppose $x_{d-1} > n-\max\{k, 0\}$. From the first step, since there exists $1 \leq i' \leq d-2$ such that $x_{i'} + \cdots + x_{d-2} = \max\{k, 0\}$, we have $x_{i'} + \cdots + x_{d-1} > n$, implying $x \notin n(\mathrm{St}_{d-1}^\ast \times [-1, 1]) \cap \ZZ^d$, which contradicts $x \in n\Ac$. Therefore, we conclude that $x_{d-1} \leq n-\max\{k, 0\}$. 

Next, suppose $x_{d-1}+x_d \leq n-\max\{k, 0\}$. From $(x_1, \ldots, x_{d-2}) \in n\mathrm{St}_{d-2}^\ast \cap \ZZ^{d-2}$, we know that $x_{i} + \cdots + x_{d-2} \leq \max\{k, 0\}$ for all $1 \leq i \leq d-2$. Consequently, $x_{i} + \cdots + x_d \leq n$ for all $1 \leq i \leq d$, implying $x \in n\mathrm{St}_d^\ast\cap \ZZ^d$, which contradicts $x \in n\mathrm{St}_d^\ast \cap \ZZ^d$. Thus, we deduce $x_{d-1}+x_d > n-\max\{k, 0\}$.
Therefore, we have $x \in \bigsqcup_{-n \leq m \leq n}(A_m \times \Delta_m) \cap \ZZ^d$.

Now, let $x = (x_1, \ldots, x_d) \in \bigsqcup_{-n \leq m \leq n}(A_m \times \Delta_m) \cap \ZZ^d$. Then there exists $-n \leq k \leq n$ such that $(x_1, \ldots, x_d) \in A_k \times \Delta_k$. Since $x_i + \cdots + x_{d-2} \leq \max\{k, 0\}$ holds for all $1 \leq i \leq d-2$ by $(x_1, \ldots, x_{d-2}) \in A_k$, we have $x_i + \cdots + x_{d-1} \leq n$. 
Additionally, since there exists $1 \leq i' \leq d-2$ such that $x_{i'} + \cdots + x_{d-2} = \max\{k, 0\}$ and $x_{d-1} + x_d > n-\max\{k, 0\}$ by $(x_{d-1}, x_d) \in \Delta_k$, we conclude $x_{i'} + \cdots + x_d > n$. Therefore, $x \in n\Ac \cap \ZZ^d$.

\smallskip

\noindent
\underline{The third step}: From (3.3) and (3.7), we have the following equation:
$$\left| \bigsqcup_{-n \leq m \leq n}(A_m \times \Delta_m) \cap \ZZ^d \right|= (2n+1)E_{\mathrm{St}_{d-1}^\ast}(n) - E_{\mathrm{St}_d^\ast}(n).$$
From this equation and (3.6), we get (3.1).
\end{proof}

\end{document}
